\documentclass{report}
\usepackage{amsmath,amssymb}
\setlength{\parindent}{0mm}
\setlength{\parskip}{1em}
\begin{document}
\begin{center}
\rule{6in}{1pt} \
{\large Misha Kilmer \\
{\bf Tensor-diamond Approximation for Construction of Kronecker Product-based Preconditioners }}

Dept of Mathematics \\ Tufts University \\ 503 Boston Ave \\ Medford \\ MA 02155
\\
{\tt misha.kilmer@tufts.edu}\\
Tamara Kolda\end{center}

\newcommand{\bfA}{{\bf A}}
\newcommand{\bfb}{{\bf b}}
\newcommand{\bfx}{{\bf x}}
\newcommand{\bfe}{{\bf e}}
\newcommand{\bfM}{{\bf M}}

\newcommand{\Mout}{\diamondsuit}


\newcommand{\M}[1]{{\bf {\boldmath {\MakeUppercase{#1}}}}}

\newcommand{\T}[1]{\boldsymbol{\cal {\MakeUppercase{#1}}}}


\newcommand{\TA}{\T{A}}


\begin{center} Tensor-diamond Approximation for Construction of Kronecker
Product-based Preconditioners \end{center}

Suppose $\bfA$ is a square, invertible, block $\ell \times \ell$ matrix
with $n \times n$ blocks, that is characterized by either (or both) of
the following
conditions:
\begin{itemize}
\item $\bfA$ is block Toeplitz (or can be permuted to be block Toeplitz);
\vspace*{-8pt}
\item $\bfA$ has small block-bandwidth $k \ll \ell$ (or, more generally,
the number of non-zero blocks is much less than $\ell$).
\end{itemize}
Such matrices might arise, for example, in image deblurring or
discretization of a PDE. In this talk, we will primarily focus on
on image deblurring applications.

We consider the problem of determining an approximation $\bfM$ to $\bfA$ such that either
\[ \bfM = \bfA_R \otimes \bfA_L \qquad \mbox{ or } \qquad \bfM = (\M{I}
\otimes \bfA_L)\M{C}, \]
where $\otimes$ denotes the Kronecker product, and $\M{C}$ is a matrix
with circulant blocks.
The leftmost problem has been studied previously (see, for example, [1,2]).
Our approach differs from previous approaches in that we produce the
desired factorization by representing
the approximation problem as a third-order tensor (i.e. 3D array) approximation problem.
In the context of PDEs, the approximation $\bfM$ has the advantage of
being relatively easy to directly invert and therefore is attractive as a
preconditioner.
In the context of image deblurring, we can relatively easily compute
rank-revealing factorizations of $\bfM$, which can then be used to
generate regularized preconditioners
or surrogate filters, etc.

Specifically,
we take the relevant descriptive information from $\bfA$ and put this
into an appropriately oriented $n \times m \times n$, third-order tensor,
denoted $\TA$.
We then consider the tensor approximation problem
\[ \min_{\M{G},\M{H}} \| \TA - \M{G} \Mout \M{H} \|_F, \qquad \mbox{with
possible constraints on \M{H}},\]
where the $\M{G} \Mout \M{H}$ is the product-cyclic tensor, first defined
in [3], given by
\[ \M{G} \Mout \M{H} \in \mathbb{R}^{n \times m \times n}, \qquad (\M{G}
\Mout \M{H})_{:,:,j} = \M{G} \M{Z}^{j-1} \M{H}^T, \]
and $\M{Z}$ is the $n \times n$ downshift matrix.
We will discuss how the entries of $\bfA_L, \bfA_R$ or $\bfA_R, \M{C}$ are
constructed from entries in $\M{G}$ and $\M{H}$.

Numerical results for spatially variant blurring operators show the
possible utility of this approach.

\vspace*{8pt}

{\bf References:} \newline
\noindent{[1]} C.~F.~{Van Loan} and N.~P.~Pitsianis, ``Approximation with
{K}ronecker Products," in {\bf Linear Algebra for Large Scale and Real
Time Applications},
Kluwer Publications, 1993, M. S. Moonen and G. H. Golub eds, pages 293-314. \newline

\noindent{[2]} J. Kamm and J. G. Nagy, ``Optimal {K}ronecker Product
Approximation of Block {T}oeplitz Matrices," SIAM Journal on Matrix
Analysis and Applications,
volume 22, 2000, pages 155-172. \newline

\noindent{[3]} M. E. Kilmer and T. G. Kolda, ``Approximations of Third Order Tensors as
Sums of (Non-negative) Product Cyclic Tensors," Plenary Presentation,
Householder Symposium, June 2011.


\end{document}
